%% 英文摘要页


\begin{center}\heiti\sanhao\textbf{
        Abstract
    }
\end{center}

With the growing demand for low-cost intelligent unmanned systems, real-time recognition and accurate aiming at fixed targets on mobile platforms has become a prominent research topic. Existing solutions are typically hindered by high cost, deployment complexity, or insufficient real-time performance. This thesis proposes and implements a fixed-target recognition and aiming system for mobile platforms, designed under the constraints of less than 1000 RMB and ``pure vision + IMU tight coupling,'' providing an engineering-practical low-cost solution.

The system adopts a three-node distributed heterogeneous architecture. The Linux vision node handles target detection, PnP pose estimation, and tracking state management. The STM32H743 gimbal controller, running the RT-Thread real-time operating system at 480\,MHz, executes a dual-loop attitude control algorithm with IMU feedforward compensation. The MSPM0G3507 chassis controller independently manages eight-sensor grayscale line tracking and differential-drive encoder PI speed control. The three nodes communicate via serial protocols and operate without mutual interference.

For visual detection and tracking, a three-stage filtering strategy is adopted: HSV color threshold segmentation, contour geometric constraint verification, and PnP reprojection size validation. This enables robust detection and center localization of A4-paper targets. A three-thread plus dual-bounded queue architecture ensures stable visual throughput above 25\,FPS.

For gimbal control, a dual-loop architecture is proposed: an outer APID loop driven by pixel error, and an inner feedforward loop driven by IMU angular velocity and acceleration. The The feedforward design is inspired by Active Disturbance Rejection Control (ADRC). Under platform acceleration or vibration disturbances, the architecture significantly reduces phase lag, achieving a reduction of more than 30\% reduction in peak aiming-line error in experiments. The steady-state pointing error satisfies the design target of $\leq$0.8\,mrad. The HiPNUC six-axis IMU provides 100\,Hz angular velocity and attitude output, decoded via the \texttt{hipnuc\_dec} protocol module.

For chassis control, a differential kinematics model decomposes forward and angular velocity commands into individual wheel speed targets. Each wheel employs an incremental PI speed control with AB-phase encoder feedback. The line-following strategy combines normalized grayscale weighted deviation estimation with a first-order low-pass filter to suppress control jitter. Additionally, an A* heuristic path planning algorithm is integrated for autonomous point-to-point maneuvering in structured environments.

Experimental results demonstrate that the visual frame rate stably exceeds 25\,FPS under standard indoor lighting, the IMU maintains 100\,Hz sampling, and the chassis completes multi-lap line-following tasks reliably. The total hardware cost remains within the budget, validating the feasibility and engineering value of the proposed system.

\textbf{KEY WORDS:} Visual Servo, Target Recognition, Gimbal Control, IMU Feedforward, Mobile Robot

\newpage\quad
\clearpage
